% The four parts of the Day 13 lab. The steps follow Labs/day13/README.md.
%
% A value that only the lab hardware can give is written as "..." and the
% students measure it.

% ------------------------------------------------------------------- Part A
\begin{frame}[fragile]{Part A: Tasks A1 and A2 (20 min)}
  \small
  \renewcommand{\arraystretch}{1.08}
  \begingroup\centering
  \begin{tabular}{@{}L{0.07\textwidth}L{0.70\textwidth}L{0.12\textwidth}@{}}
    \toprule
    \textbf{Step} & \textbf{Work} & \textbf{Time} \\
    \midrule
    1 & Copy the lab folder, check the Day 8 folder and its model & 5 min \\
    2 & Tasks A1 and A2 in \code{pi/monitor.py} & 15 min \\
    \bottomrule
  \end{tabular}\par\endgroup
\begin{lstlisting}[style=shell]
scp -r day13 edge@pi-07.local:edgeai/                           # laptop
ls ~/edgeai/day08/pi/detector.py ~/edgeai/day08/models/cupbottle_320_ncnn_model
cd ~/edgeai/day13/pi && ~/yolo/bin/python monitor.py            # Pi
\end{lstlisting}
  \begin{itemize}
    \item \textbf{Task A1} \code{input\_stats(frame)}: the grey image with
          \code{cv2.cvtColor}, its mean (brightness), and the variance of
          \code{cv2.Laplacian} (sharpness).
    \item \textbf{Task A2} \code{summarize(records, seconds)}: the ten keys
          of the docstring. The p95 uses \code{percentile()}, the
          nearest-rank rule of Day 9.
    \item \textbf{You see:} \code{Task A1: complete} and \code{Task A2:
          complete}.
  \end{itemize}
\end{frame}

\begin{frame}[fragile]{Part A: run the detector, read the log (20 min)}
  \small
  Put the cup and the bottle in front of the camera.
\begin{lstlisting}[style=shell]
~/yolo/bin/python monitor_detect.py \
    --model ../../day08/models/cupbottle_320_ncnn_model \
    --group g07 --no-mqtt --seconds 120
wc -l -c ../logs/detector.jsonl
tail -n 1 ../logs/detector.jsonl | python3 -m json.tool
\end{lstlisting}
  \textbf{You see} one line each 10 s (here in two parts):
\begin{lstlisting}[basicstyle=\ttfamily\tiny]
10:21:40   ... frames/s  latency .../... ms  confidence ...  low ...  objects ...
  brightness ...  sharpness ...  ... C  rss ... MB
\end{lstlisting}
  \begin{itemize}
    \item The log has one JSON record on each line: \code{start}, one
          \code{summary} each 10 s, \code{stop}. It has at most 1 MB and 5
          old files.
    \item Cover the lens for 10 s: which values change?
    \item \textbf{You record:} one line, the records and the bytes of the
          log, and the WARNING records of a run with \code{--slow-ms 50}
          (\code{grep -c slow\_frame}).
  \end{itemize}
  \source{The values of the lines come from your board.}
\end{frame}

% ------------------------------------------------------------------- Part B
\begin{frame}[fragile]{Part B: the live dashboard (25 min)}
  \small
\begin{lstlisting}[style=shell]
systemctl is-active mosquitto                       # active (Day 12)
cd ~/edgeai/day13/dashboard
~/yolo/bin/python dashboard.py --group g07          # second SSH terminal
cd ~/edgeai/day13/pi                                # first terminal
~/yolo/bin/python monitor_detect.py \
    --model ../../day08/models/cupbottle_320_ncnn_model --group g07
\end{lstlisting}
  \begin{itemize}
    \item Open \code{http://pi-07.local:8080} on the laptop. The state bar
          says \code{Waiting for data ...}, then \code{no alert program}.
    \item The eight charts get one point each 10 s: frame rate, latency
          p95, temperature, memory of the detector, mean top score, part
          below 0.5, objects, brightness.
    \item On the laptop:
          \code{mosquitto\_sub -h pi-07.local -t 'edgeai/g07/pi/\#' -v}.
    \item \textbf{You record:} the messages and the lost messages of the
          state bar after 5 minutes, and the name of the CSV file in
          \code{dashboard/metrics\_log/}.
  \end{itemize}
\end{frame}

\begin{frame}{Part B: reference period and Task C1 (20 min)}
  \begingroup\centering
  \begin{tikzpicture}[font=\scriptsize, >=stealth, x=0.85cm]
    \draw[->, coolgray] (0,0) -- (12.4,0) node[below left] {time};
    \fill[treegreen!35] (0.2,0.12) rectangle (6.0,0.82);
    \node at (3.1,0.47) {reference: 10 min of the normal scene};
    \fill[darkorange!40] (6.6,0.12) rectangle (8.6,0.82);
    \node[align=center] at (7.6,0.47) {dim light\\ 3 min};
    \fill[darkorange!25] (9.6,0.12) rectangle (11.6,0.82);
    \node[align=center] at (10.6,0.47) {unknown\\ object};
    \node[below] at (0.2,0) {B4}; \node[below] at (6.6,0) {C2};
  \end{tikzpicture}\par\endgroup
  \vskip 0.5em
  \small
  \begin{itemize}
    \item Keep the scene normal for 10 minutes: the cup and the bottle in
          view, the normal light. People can move as usual. Do not stop the
          detector.
    \item Write the start and the end time. Write which panels move and
          which stay flat.
    \item During this time, complete \textbf{Task C1}
          \code{threshold\_from\_history(values, window, k)} in
          \code{pi/reference.py}: the means of 6 summaries in a row, their
          mean and standard deviation, the threshold mean $-$ 3 sd, and the
          clear value mean $-$ 1 sd.
  \end{itemize}
  \keypoint{The reference must show the normal variation of your scene.
    A still scene with no variation gives a threshold that fires too
    often.}
\end{frame}

% ------------------------------------------------------------------- Part C
\begin{frame}[fragile]{Part C: a threshold from the reference (10 min)}
\begin{lstlisting}[style=shell]
~/yolo/bin/python reference.py \
    --csv ../dashboard/metrics_log/metrics_2026-10-05.csv \
    --start 10:20 --end 10:30 --group g07
\end{lstlisting}
  \textbf{You see:}
\begin{lstlisting}[basicstyle=\ttfamily\scriptsize]
Task C1: complete
Reference: ... rows from 10:20:00 to 10:30:00, metric confidence
Rows:      mean ...  sd ...  min ...  max ...
Means of 6 rows: ... means, mean ...  sd ...
Threshold: ...   clear: ...   (mean - ..., mean - ...)
Command: python alert.py --group g07 --metric confidence --threshold ...
\end{lstlisting}
  \small
  \begin{itemize}
    \item Use the date of today and the times of your reference.
    \item A \code{WARNING: the reference has almost no variation}: run it
          again with \code{--min-margin 0.05}.
    \item Run it also with \code{--metric brightness}. \textbf{You record:}
          the mean, the standard deviation, the threshold, and the clear
          value.
  \end{itemize}
\end{frame}

\begin{frame}{Part C: two drift events (30 min)}
  \small
  Make each event for 3 minutes, with 3 minutes of the normal scene between
  them. Write the start and the end time.
  \vskip 0.4em
  \footnotesize
  \renewcommand{\arraystretch}{1.15}
  \setlength{\tabcolsep}{4pt}
  \begingroup\centering
  \begin{tabular}{@{}L{0.20\textwidth}L{0.38\textwidth}L{0.34\textwidth}@{}}
    \toprule
    \textbf{Event} & \textbf{What you do} & \textbf{Watch} \\
    \midrule
    Dim light & Switch off the light of the room, or turn the lamp away &
      Mean top score, objects, brightness \\
    Unknown object & Put your unknown object in place of the cup & Mean top
      score, part below 0.5, objects \\
    \bottomrule
  \end{tabular}\par\endgroup
  \vskip 0.6em
  \small
  \begin{itemize}
    \item \textbf{You record:} for each event, the mean top score, the part
          below 0.5, the objects, the brightness, and the sharpness of the
          printed lines (the table of the report).
    \item Which metric shows each event most clearly? Which metric names
          the cause? The model gives no error in either event.
    \item With no lamp to switch: \code{--dim-after 60 --dim-gain 0.1}
          multiplies each frame by 0.1 for 3 minutes.
  \end{itemize}
\end{frame}

% ------------------------------------------------------------------- Part D
\begin{frame}{Part D: Task D1, the alert rule (10 min)}
  \begingroup\centering
  \begin{tikzpicture}[font=\scriptsize, >=stealth, x=0.7cm, y=1.6cm]
    \draw[->, coolgray] (0,0) -- (13.6,0) node[below left] {time};
    \draw[->, coolgray] (0,0) -- (0,1.15) node[above right] {mean of 6};
    \draw[cardinalred, dashed] (0,0.45) -- (13.3,0.45);
    \node[cardinalred, above right] at (11.2,0.36) {threshold};
    \draw[treegreen, dashed] (0,0.62) -- (13.3,0.62);
    \node[treegreen, above right] at (11.2,0.6) {clear};
    \draw[darkcyan, thick] (0,0.85) -- (3,0.85) -- (5,0.3) -- (9,0.3)
      -- (10.5,0.8) -- (13,0.85);
    \fill[cardinalred!20] (4.3,0.02) rectangle (6.8,0.12);
    \node[above, text=cardinalred] at (5.55,0.1) {duration};
    \draw[->, cardinalred] (6.8,0.95) node[above] {fire} -- (6.8,0.3);
    \draw[->, treegreen] (10.1,0.95) node[above] {clear} -- (10.1,0.66);
  \end{tikzpicture}\par\endgroup
  \vskip 0.4em
  \small
  \begin{itemize}
    \item \textbf{Task D1} \code{AlertRule.update(t, value)} in
          \code{pi/alert.py}: append the value; the mean of the last 6
          values (60 s); fire when the mean stays below the threshold for
          the duration; clear when the mean is above the clear value.
    \item The clear value is above the threshold (hysteresis), so a mean
          near the threshold does not switch the alert on and off.
    \item The docstring gives the five steps. The program checks the
          method with a known sequence before it connects.
  \end{itemize}
\end{frame}

\begin{frame}[fragile]{Part D: test the alert (15 min)}
\begin{lstlisting}[style=shell]
~/yolo/bin/python alert.py --group g07 --metric confidence \
    --threshold ... --clear ... --window 6          # the command of C1
\end{lstlisting}
  \textbf{You see} during the dim light event:
\begin{lstlisting}[basicstyle=\ttfamily\scriptsize]
Task D1: complete
Rule: mean of 6 summaries of confidence below ... for 60 s; clear above ...
HH:MM:SS ALERT ON    confidence: mean of 60 s ..., threshold ..., clear ...
HH:MM:SS ALERT OFF   confidence: mean of 60 s ..., threshold ..., clear ...
\end{lstlisting}
  \small
  \begin{itemize}
    \item The state bar of the dashboard is red during the alert. With the
          XIAO of Day 12, \code{--xiao-led} switches its LED.
    \item A second copy with \code{--metric brightness} and its own
          threshold can watch the light.
    \item \textbf{You record:} the seconds from the start of the event to
          \code{ALERT ON}, and from its end to \code{ALERT OFF}. Then write
          the incident report.
  \end{itemize}
\end{frame}
